\begin{resumo}
O Projeto EMA, conduzido no Laboratório de Automação e Robótica (LARA) da Universidade de
Brasília, desenvolve sistemas de reabilitação motora que integram eletroestimulação
funcional, robótica e realidade virtual. Ao longo de anos, sucessivas gerações
de estudantes produziram diversos subsistemas que, embora funcionais do ponto de vista da
pesquisa, apresentam fragilidades recorrentes de Engenharia de Software: arquiteturas
acopladas, ausência de documentação, versionamento desorganizado e uma lacuna sistemática
entre o que está registrado e o que está de fato montado. Este trabalho, de natureza
aplicada e exploratória, aplica o instrumental da Engenharia de Software ao diagnóstico,
à documentação e à proposta de reestruturação desses sistemas, adotando uma postura de
consultoria: assume-se a perspectiva de um novo integrante que tenta reproduzir os sistemas
tal como se encontram, registrando cada obstáculo como dado de diagnóstico. A investigação
organiza-se em dois estudos de caso complementares e em estágios opostos de maturidade: o
sistema de realidade virtual imersivo, inativo há anos, cujo desafio é a
reativação; e o sistema de ciclismo assistido por eletroestimulação funcional, ativo, porém
de difícil reprodução. Para cada caso, documenta-se a arquitetura atual, relata-se a
tentativa de reprodução e mapeiam-se as limitações específicas; em seguida, consolidam-se os
problemas comuns, com destaque para a organização do versionamento e a documentação. Como
resultado, o trabalho entrega o diagnóstico estruturado, a diagramação padronizada das
arquiteturas e um conjunto de recomendações de reestruturação, preparando o terreno para a
etapa de integração a ser conduzida em trabalho posterior.

\vspace{\onelineskip}

\noindent
\textbf{Palavras-chave}: engenharia de software. robótica de reabilitação. documentação. refatoração. projeto EMA.
\end{resumo}
